\begin{definition} \label{definition:ample_divisor_on_a_scheme}
Let $X$ be a \CrefAndHyperrefIfExist{definition:noetherian_scheme}{Noetherian scheme}. Let $\mathcal{L}$ be an invertible sheaf (line bundle) on $X$.

The invertible sheaf $\mathcal{L}$ is said to be \hldef{ample} if for every coherent sheaf $\mathcal{F}$ on $X$, there exists an integer $n_0 \geq 0$ such that for all $n \geq n_0$, the twisted sheaf $\mathcal{F} \otimes_{\mathcal{O}_X} \mathcal{L}^{\otimes n}$ is generated by its global sections. That is, the evaluation morphism
$$\hlin{\bigoplus_{s \in \Gamma(X, \mathcal{F} \otimes \mathcal{L}^{\otimes n})} \mathcal{O}_X \to \mathcal{F} \otimes \mathcal{L}^{\otimes n}}$$
is surjective as a morphism of $\mathcal{O}_X$-modules.

An invertible sheaf $\mathcal{L}$ is said to be \hldef{very ample} (relative to a base scheme $S$, when $X$ is an $S$-scheme) if there exists a closed immersion $i: X \hookrightarrow \mathbb{P}^n_S$ for some $n \geq 0$ such that $\mathcal{L} \cong i^* \mathcal{O}_{\mathbb{P}^n_S}(1)$, where $\mathcal{O}_{\mathbb{P}^n_S}(1)$ is the Serre twisting sheaf (\Cref{definition:twisting_sheaf_of_serre_of_degree_n_on_projective_space_over_a_scheme}).

A Cartier divisor $D$ on $X$ is said to be an \hldef{ample divisor} if the associated invertible sheaf $\mathcal{O}_X(D)$ is ample. Similarly, $D$ is called a \hldef{very ample divisor} if $\mathcal{O}_X(D)$ is very ample.

\TextIfExists{definition:projective_morphism_of_schemes_projective_scheme_over_a_scheme}{
When $X$ is a scheme over a base scheme $S$, an invertible sheaf $\mathcal{L}$ on $X$ is said to be \hldef{$S$-ample} (or relatively ample) if for every coherent sheaf $\mathcal{F}$ on $X$, there exists an integer $n_0 \geq 0$ such that for all $n \geq n_0$, the sheaf $\mathcal{F} \otimes_{\mathcal{O}_X} \mathcal{L}^{\otimes n}$ is generated by global sections relative to $S$. That is, the natural map
$$f^* f_* (\mathcal{F} \otimes \mathcal{L}^{\otimes n}) \to \mathcal{F} \otimes \mathcal{L}^{\otimes n}$$
is surjective, where $f: X \to S$ is the structure morphism.

A Cartier divisor $D$ on $X$ is called an \hldef{$S$-ample divisor} if $\mathcal{O}_X(D)$ is relatively ample over $S$. By Serre's Criterion for projectivity, a \CrefAndHyperrefIfExist{definition:proper_morphism_of_schemes}{proper} morphism $f: X \to S$ of Noetherian schemes is projective if and only if there exists an $S$-ample invertible sheaf on $X$.
}

The following equivalent characterizations of ampleness hold for a Noetherian scheme $X$:
\begin{enumerate}
    \item $\mathcal{L}$ is ample if and only if for every coherent sheaf $\mathcal{F}$, there exists an integer $n > 0$ such that $\mathcal{F} \otimes \mathcal{L}^{\otimes n}$ is generated by global sections.
    \item $\mathcal{L}$ is ample if and only if the open subsets $X_f = \{x \in X : f(x) \neq 0\}$, where $f$ ranges over all elements of $\Gamma(X, \mathcal{L}^{\otimes n})$ for all $n > 0$, form a basis for the topology of $X$.
    \item If $X = \operatorname{Spec}(A)$ is an affine Noetherian scheme, then the structure sheaf $\mathcal{O}_X$ is ample.
\end{enumerate}

When $S = \operatorname{Spec}(k)$ for a field $k$, and $X$ is a proper scheme over $k$, an invertible sheaf $\mathcal{L}$ is ample if and only if for every coherent sheaf $\mathcal{F}$, the cohomology groups $H^i(X, \mathcal{F} \otimes \mathcal{L}^{\otimes n})$ vanish for all $i > 0$ and all sufficiently large $n$.
\end{definition}
